\section{Proactive：AI 产品和互联网产品的本质区别}

Tristan 认为，AI 时代最大的交互变化不是 LUI 或 GUI，而是 proactive：终于有东西能主动替你做事。OpenClaw 让大家在工具层感受到 proactive，Elys 则把 proactive 放到社交层，让分身主动面对世界。

\begin{figure}[H]
\centering
\includegraphics[width=0.92\textwidth]{proactive-product-shift.png}
\caption{从 Reactive 到 Assistive 再到 Proactive 的产品范式。}
\end{figure}

\begin{knowledgebox}{读图：Proactive 改变了产品心智}
Reactive 产品需要用户发起；Assistive 产品在用户操作时辅助；Proactive 产品能在用户不主动操作时替用户探索机会。Elys 的想象力来自第三层：分身主动社交并带回连接。
\end{knowledgebox}

\subsection{和模型公司的竞争}

Tristan 认为，AI 社交不一定需要自己训练基础模型。原因是社交网络的核心不只是智能，而是 context、产品范式和真人连接。模型公司关注智能本身，未必把社交作为第一优先级；应用公司若能掌握 context 飞轮，也可能有独立机会。

\begin{importantbox}{应用公司的壁垒在哪里}
不是“我也能调用模型”，而是我是否拥有独特 context、真实关系、产品心智、数据反馈和用户信任。没有这些，应用会被模型公司覆盖；有这些，应用可能形成自己的网络效应。
\end{importantbox}

\subsection{人类太孤独了}

访谈最后回到精神内核：人类太孤独了。Eve 试图通过 AI 陪伴解决孤独，Elys 试图通过 AI 促进真人连接解决孤独。前者是 AI 与人的关系，后者是 AI 帮人和人建立关系。这个区别决定了 Elys 不是陪伴产品，而是社交网络产品。

\begin{warningbox}{情绪价值不是低级需求}
社交、陪伴、被理解和连接是高频而深层的人类需求。AI 产品若能安全、真实、可控地改善这些需求，可能比纯工具产品更接近普通用户。但风险也更高，因为它触及隐私、身份和关系。\end{warningbox}

\subsection{本章小结}

Proactive 是 AI 产品区别于互联网产品的关键。Elys 的核心不是让 AI 陪聊，而是让 AI 分身主动替用户降低社交摩擦，最终带回真人连接。


